%!TEX root = ../main.tex
\chapter{Experimental Setup and Evaluation Metrics}
\label{ch:experiments}

Both training regimes used in this work---imitation of recorded flights
(Section~\ref{sec:training-supervised}) and reinforcement learning in
simulation (Section~\ref{sec:training-rl})---produce a single headline number
that the optimiser minimises: a loss in the first case, a negative expected
return in the second. Neither number answers the question that actually
matters, which is whether the resulting network flies the robot well. A
supervised loss can fall while the network learns nothing but ``predict zero'';
an episodic return can rise while the vehicle learns to crash less rather than
to reach its target. This chapter fixes, before any result is shown, the
conditions under which the numbers of Chapter~\ref{ch:results} were produced:
the machines and the software the runs were executed on, the protocol every run
follows, and every quantity that is logged alongside the headline number, what
each one measures and why it was added. The metrics are grouped by training
regime, since the two ask different questions of the model, and each group ends
with a summary table.

Every metric named here is a column in a CSV file written during training: one
row per epoch for the supervised trainer, one row per iteration for the
reinforcement trainer. The column names are quoted verbatim in monospace, so
that a figure in Chapter~\ref{ch:results} can always be traced back to its
definition here.

The chapter is organised in two halves. The first fixes the common ground: the
notation used throughout (Section~\ref{sec:metrics-notation}), the
infrastructure the experiments run on (Section~\ref{sec:exp-tooling}) and the
protocol under which the architectures are trained, validated and compared
(Section~\ref{sec:exp-protocol}). The second defines the two families of
metrics, one per training stage: supervised in
Section~\ref{sec:metrics-supervised}, reinforcement learning in
Section~\ref{sec:metrics-rl}.

\subsection{Conventions and Notation}
\label{sec:metrics-notation}

The network's output is a velocity command
$\mathbf{y} = (v_x, v_y, v_z, \omega_x, \omega_y, \omega_z) \in \mathbb{R}^{6}$:
three linear velocities in \si{\metre\per\second} and three angular velocities
in \si{\radian\per\second}. Components are indexed $i = 0,\dots,5$ in the order
above, matching the \met{_dim_<i>} suffix of the log columns. The demonstrated
(target) command is $\mathbf{y}$ and the network's prediction is
$\hat{\mathbf{y}}$; $N$ denotes the number of samples over which an average is
taken, which in sequence mode counts every time step of every window.

Each robot in the fleet declares its own kinematic limits inside the
\texttt{state} block of the observation. The per-component \emph{velocity cap}
$c_i > 0$ is the largest magnitude the platform accepts on component $i$;
a cap of zero declares that the axis does not exist on that robot (for example
vertical velocity on a ground vehicle).

For the reinforcement learning metrics the following terms are used with their
usual meaning, defined formally in Section~\ref{sec:rppo-gae}. An
\emph{episode} is one attempt at the navigation task, from spawn to a terminal
condition or to the time limit. A \emph{rollout} is a fixed number of
environment steps collected from every parallel environment with the current
policy; an \emph{iteration} is one rollout followed by one policy update, and
a \emph{mini-batch} is a subset of the rollout the update is computed on.
$\pi_\theta(a \mid s)$ is the stochastic policy, $V(s)$ the critic's estimate of
the state value, $\gamma$ the discount factor, $A_t$ the advantage of the action
taken at step $t$ and $R_t$ the corresponding bootstrapped return.

\subsection{Infrastructure and Tooling}
\label{sec:exp-tooling}

\subsubsection{Hardware and Software Environment}

Training is carried out on a single workstation equipped with an NVIDIA RTX~5080
with \SI{16}{\gibi\byte} of video memory and \SI{32}{\gibi\byte} of system
memory, on which both the training procedures and the fleet of simulated
environments run. The memory of the accelerator is the resource that constrains
the experimental design more than any other: it is shared between the training
session and the slices of the fleet, and it determines the finest grid on which
the encoder may be refined during the reinforcement stage
(Section~\ref{sec:rl-networks}).

The software environment is the one declared by the two repositories. The
networks and the training are implemented in PyTorch with CUDA support; the
sparse convolution uses \texttt{spconv}; the corpus is read as a stream through
\texttt{webdataset}; the control environment comes from \texttt{gymnasium}. The
pipeline of Chapter~\ref{ch:pipeline} and the simulators run instead in Docker
containers, so that the version of ROS~2, of the simulator and of the perception
libraries is fixed by the image rather than by the machine running it.

\subsubsection{Simulation and the Parallel Fleet}

The reinforcement stage requires many environments in parallel, and each of them
is an isolated copy of the entire system rather than a simplified model
(Section~\ref{sec:rl-environment}). The fleet is managed by a dedicated tool
that governs its life cycle in five operations: \emph{planning}, which works out
how many slices to instantiate and with which configurations, recording every
runtime parameter in a manifest; \emph{materialisation}, which generates the
files of each slice from the manifest; \emph{start-up}, which brings up the
containers and the processes in the required order; \emph{inspection}, which
reports what is running on each slice and with which configuration; and
\emph{shutdown}.

Two properties of this tool belong to the experimental protocol rather than to
the infrastructure. The first is that start-up checks which slices are actually
producing point clouds and writes a reduced manifest of them: starting a session
on a manifest that lists a slice which is not working means losing the session a
few minutes after beginning it. The second is that the number of slices can be
derived from the machine, from a measured per-slice cost and a reserve for the
training session, which runs on the same machine and itself requires cores,
memory and accelerator memory. The number so obtained states how many slices
\emph{fit} on the machine, not how many it is worth running: beyond a certain
number the overall throughput grows less than proportionally, and the choice
made in the experiments is reported in Chapter~\ref{ch:results}.

\subsubsection{Logging and Run Artefacts}

Every session produces artefacts in two distinct places, and the separation is
deliberate: what is needed to \emph{read} a session is kept apart from what is
needed to \emph{resume} it or to transfer it to a robot.

The supervised stage writes one row per epoch to a comma-separated file,
containing every metric defined in Sections~\ref{sec:metrics-supervised}
and~\ref{sec:metrics-rl}, plus a console summary; a finer per-step record can be
enabled when an epoch is long enough to make the epoch average uninformative
about what happened inside it. The reinforcement stage likewise writes one row
per iteration and a textual log. The checkpoints, described in
Section~\ref{sec:module-implementation}, live in a separate tree, one directory
per save.

A property on which Chapter~\ref{ch:results} relies follows: every figure
reported there is generated from one of those files, and every column of those
files is defined in this chapter. No number reported in what follows was
transcribed by hand from a console.

\subsubsection{Plotting and Analysis Tools}

The figures are produced by two generators, one per stage, which read the log
files and derive from them one figure per family of metrics, plus a summary
document. Both accept several sessions at once and overlay them, and both omit
the figures for which a session did not record the corresponding columns --- so
that comparing a session with the auxiliary term against one without it requires
no configuration. The style is shared: a validated palette in a fixed order, a
distinct marker per series so that a curve is never identified by colour alone,
and labelled reference lines rather than legend entries.

Beside them stands a set of measurement tools that do not belong to the training
and each of which answers a stated question. Two concern the encoder and are
used in Chapter~\ref{ch:results}: \emph{ablation}, which replaces the encoder's
output with the batch mean and measures how much the loss degrades, answering
the question of whether the command uses the map; and the \emph{probe}, which
fits a linear read-out of the frozen descriptor towards the observable geometry,
answering the distinct question of whether the encoder represents the scene. The
latter can also be evaluated on the same architecture untrained, which gives the
floor a trained read-out must clear in order to be called learnt.

The remaining tools concern the fleet and the reward: the \emph{pricing} of a
reward configuration in advance, discussed in Section~\ref{sec:rl-reward}; the
\emph{comparison} between the observation the fleet delivers and the one the
corpus recorded, feature by feature; the \emph{replay} of the perception signals
over the recorded corpus, which is how a terminal threshold is calibrated before
being armed; and the analysis of the final state of the episodes, which
separates causes of termination that the aggregate counters do not.

\subsection{Evaluation Protocol}
\label{sec:exp-protocol}

\subsubsection{Staged Validation: From Synthetic Tasks to the Fleet}

Neither of the two training stages was applied to the problem before being
verified on tasks whose answer is known. The progression has three levels, and
the reason it exists is that on this problem an implementation defect produces
not an error but a plausible curve.

The first level consists of automated checks that run without the fleet, without
clouds and without an accelerator, on synthetic data. They measure no
performance: they pin down the properties the training depends on and that a
change might silently violate --- that the two segments of a flight fall on the
same side of the split, that a mixed batch contains a single grid, that
validation is not reflected, that the reflection is its own inverse and swaps
every lateral quantity, that the numbering of the sectors is the same in the
training and in the measuring tool, and that a policy update moves what it must
and leaves still what it must not.

The second level is a synthetic task for the supervised stage: next-step
prediction of a sine wave with randomly drawn phase and frequency and additive
noise. It is deliberately trivial as a problem and informative as a check,
because it exercises the whole recurrent path --- windows, state propagation,
truncated backpropagation --- on a task whose solution is known, so that a cell
which does not learn there is not to be looked for in the corpus.

The third level is a standard control environment for the reinforcement stage,
with a dense observation and no fleet. It exercises the algorithm in full ---
collection, advantage estimation, trust region, normalisers, recurrence --- in
an environment whose expected behaviour is known, and it is by construction the
simplest case of the interface: a change to the agent that breaks this one
breaks every environment. Two limits must however be stated, because
Chapter~\ref{ch:results} must not credit it with more than it verifies. Its
action space is discrete, so that it does not exercise the Gaussian policy of
Section~\ref{sec:rl-actions}, its learnt width or the envelope penalty. And its
time step is constant, so that it does not exercise the handling of the variable
interval which is the reason continuous-time cells were chosen.

Only after these three levels are the stages applied, respectively, to the
recorded corpus and to the fleet. Two checks precede every session on the fleet:
verifying that the manifest describes the heterogeneous fleet intended and not
several copies of the same configuration
(Section~\ref{sec:rl-environment}), and pricing the reward configuration in
advance (Section~\ref{sec:rl-reward}).

\subsubsection{Data Splits and Held-Out Episodes}

The corpus is divided into three partitions --- training, validation and test
--- according to the run-stratified criterion of Section~\ref{sec:sl-split}. The
fractions are fixed once and kept across all the sessions compared, and the seed
governing them is distinct from the session's, so that repeating an experiment
does not change the data on which it is evaluated.

The validation partition is not a neutral set, and the statement must be made
explicitly because it conditions the reading of the results: on it the
coefficients of the output calibration are estimated
(Section~\ref{sec:sl-calibration}) and the checkpoint is selected
(Section~\ref{sec:sl-optimisation}). It is therefore used twice to decide, and
the quantities measured on it do not constitute an unbiased estimate of
generalisation. That is precisely why the third partition exists, which takes
part in no decision and is read once, at the end, on the selected models alone.

\subsubsection{Checkpoint Selection and Stopping Criteria}

For the supervised stage every epoch produces a checkpoint and no automatic
interruption takes place (Section~\ref{sec:sl-optimisation}). The checkpoint
selected is the one that minimises the validation loss, but the selection is not
read on that number alone: the loss is compared against the reference predictors
of Section~\ref{sec:metrics-r2} and against the amplitude metrics of
Section~\ref{sec:metrics-amplitude}, since a model may improve the loss by
reducing what it commands, which is exactly the behaviour a squared error
rewards. An epoch that improves the loss while degrading those quantities is
flagged in Chapter~\ref{ch:results} rather than selected in silence.

For the reinforcement stage there is no single number defining the best
checkpoint, since the return depends on the reward configuration and is not
comparable across different configurations. Checkpoints are saved at a fixed
interval of iterations, and the selection is made on the task metrics of
Section~\ref{sec:metrics-rl-task}, which describe the behaviour rather than its
valuation. What is stated in Section~\ref{sec:rl-finetuning} also applies: the
initial iterations reserved for the critic constitute the measurement of the
imitated policy on the fleet, and it is with respect to that measurement --- and
not with respect to zero --- that the gain of the refinement is to be read.

\subsubsection{Comparison Across Architectures}

The comparison between the six architectures is conducted holding equal
everything that is not the cell and the wiring: the same observation and the
same normalisation, the same encoder and the same configuration, the same fusion,
state and head widths, the same output stage, and every training setting ---
including the corpus, the split, the seed governing it and the learning-rate
schedule. The list is given in Section~\ref{sec:networks-implementation}
together with what is instead not held equal, namely the parameter count and the
way the elapsed interval reaches the cell, and both asymmetries are to be read
alongside the result.

For the reinforcement stage the comparison is conducted at equal step budget,
equal fleet manifest and equal reward configuration, and starting from the
imitation checkpoints of the same architecture: comparing two architectures from
weights of different quality would measure the imitation rather than the
refinement.

The comparison is read neither on the loss alone nor on the return alone, for
the reason set out at the opening of this chapter, but on the families of
metrics of Sections~\ref{sec:metrics-supervised} and~\ref{sec:metrics-rl}, which
were introduced precisely in order to separate what those two numbers conflate.

\subsubsection{Repetitions, Seeds and Reported Variability}

A difference between two configurations is meaningful only if it exceeds the
difference between two runs of the same configuration. The protocol therefore
adopts the following rule: where two configurations are compared, the same
configuration is run more than once changing the initialisation seed alone, and
the spread between those runs is the threshold a difference must beat in order
to be stated. Since the split seed stays fixed, that spread measures the
variability of the session and not that of the data.

Two caveats complete the rule, and are stated here because
Chapter~\ref{ch:results} abides by them. The first is that on the fleet the
variability also includes the asynchrony of the slices --- real simulators
running in separate containers, whose interleaving depends on the load of the
machine, so that two identically configured runs do not collect the same
sequence of transitions --- which a single run does not allow to be separated
from the effect one intends to measure. A run on the fleet is therefore not
reproducible bit for bit: what the seed fixes is the task, that is which
episodes, from which positions and on which slices, and not the order in which
they occur. The second caveat is that the cost of a session does not allow every
configuration to be repeated: where one arm of a comparison was run only once,
the result is reported as indicative rather than as established, and the
distinction is kept explicit in the text rather than left to the reader.

\subsection{Supervised Training Metrics}
\label{sec:metrics-supervised}

The supervised trainer learns to reproduce, at each instant of a recorded
flight, the command the human pilot or the autonomous pilot gave from the observation the robot had.
The questions its metrics answer are, in order: is the model learning anything
beyond a trivial rule; \emph{what} kind of behaviour is it reproducing; and is
the machinery (encoder, recurrent cells, gradients) healthy while it does so.
Every metric below is computed per batch and averaged over the epoch, and is
reported for the training split (prefix \met{train_}) and for the validation
split (prefix \met{val_}). A third family, \met{val_cal_}, is described in
Section~\ref{sec:metrics-calibrated}.

\subsubsection{The Training Loss: Cap-Scaled Mean Squared Error}
\label{sec:metrics-loss}

The loss that is optimised, and the quantity reported as \met{train_loss} and
\met{val_loss}, is a mean squared error in which every component of the command
is first divided by the robot's declared velocity cap:
%
\begin{equation}
  \mathcal{L}
  = \frac{1}{N}\sum_{n=1}^{N}\;\frac{1}{6}\sum_{i=0}^{5}
    \left(\frac{\hat{y}_{n,i} - y_{n,i}}{s_{n,i}}\right)^{2},
  \qquad
  s_{n,i} =
  \begin{cases}
    c_{n,i} & \text{if } c_{n,i} > \varepsilon,\\
    1       & \text{otherwise},
  \end{cases}
  \label{eq:capscaled}
\end{equation}
%
with $\varepsilon = 10^{-3}$ and $c_{n,i}$ the cap declared by the robot that
recorded sample $n$.

Why the scaling is the robot's own cap rather than the standard deviation of
the corpus is argued in Section~\ref{sec:sl-loss}. What matters for reading the
number is the consequence: the error is a \emph{fraction of what that robot can
do}, so that platforms with very different limits are directly comparable and a
heterogeneous fleet needs no special handling, every sample carrying its own
caps. A cap below $\varepsilon$ is treated as ``this axis does not exist'':
the component is left in raw units, its target is necessarily zero, and the
loss simply drives the prediction to zero there without inflating its weight.

Two auxiliary values accompany the loss. The learning rate in effect at the
end of the epoch is logged as \met{lr}, so that a change of slope in the loss
curve can be matched to the schedule. When step-level logging is enabled,
\met{quick_val_loss} is the same loss evaluated on a handful of validation
batches part-way through an epoch, a cheap generalisation check for epochs that
last hours.

\subsubsection{The Families of Diagnostics}
\label{sec:metrics-families}

The loss is the quantity optimised, and on this corpus it is not by itself
evidence that anything was learnt. Six further families are logged alongside it,
each answering a question the loss cannot. They are summarised here; their
columns are defined one by one in Appendix~\ref{app:metrics}.

\paragraph{Reference predictors and skill scores.}
A loss value on its own has no scale, since it depends on how large the commands
are. Two zero-parameter predictors are therefore evaluated on every batch and
the model's error expressed relative to them: predicting zero everywhere, and
echoing the velocity the observation already contains. The first yields a
coefficient of determination, the second a stricter \emph{skill} score which is
positive only when the model earns its parameters, since a velocity command is
close to a persistent quantity and repeating it is a strong rule that costs
nothing.

\paragraph{Amplitude and restart.}
A squared error cannot see the one failure that matters most for a deployed
controller. A model trained by least squares regresses towards the conditional
mean of the demonstration: where the pilot sometimes commanded a strong
manoeuvre and sometimes none, the model learns to command a weak one every time
--- which the loss rewards and the vehicle pays for. The amplitude ratio
compares the spread of the prediction with that of the demonstration, and the
restart ratio does the same on the frames where the robot is at rest and the
demonstration starts moving, which are the rarest and the most consequential.

\paragraph{Still and moving frames.}
A substantial share of the corpus is made of legitimate ``hold still'' targets,
and they are also the easiest samples to score well on. Splitting the error by
whether the target is exactly zero on every component tells how much of a score
comes from recognising a stop and how much from reproducing a manoeuvre.

\paragraph{Per-robot breakdown.}
The corpus is recorded on more than one platform, and a fleet average can hide
one robot being served well and another badly --- the more so because a
component that is merely rare on one robot is unreachable on another. Every
ratio is therefore also reported per robot model.

\paragraph{Calibrated metrics.}
The regression to the mean described above is corrected before deployment by the
output calibration of Section~\ref{sec:output-calibration}. Since the policy that
flies is the calibrated one, every validation metric is logged twice: once on
the raw output, which is what the optimiser sees, and once through the fitted
calibration, which is what the robot executes.

\paragraph{Auxiliary geometry.}
When the auxiliary task of Section~\ref{sec:geometry-head} is enabled it
contributes its own columns, scored against the geometry each frame's own cloud
shows. They describe what the encoder has learnt about the scene, which the
imitation loss does not ask of it.

\paragraph{Model health.}
The families above describe the output. A last one describes whether the parts
of the network that produce it are doing anything at all --- the spread of the
encoder's descriptors across samples, the norms of the hidden states, the norms
of the gradients reaching each block, and the cost of an epoch. None of these is
visible on a loss curve, and each of them has at some point explained a run that
would not learn.

\subsubsection{Summary}

Table~\ref{tab:metrics-supervised} lists the supervised metrics by the
question they answer. Each is logged for the training and validation splits,
and the validation ones additionally through the calibration.

\begin{table}[htbp]
  \centering
  \small
  \caption{Supervised training metrics, grouped by the question they answer.
  The \texttt{train\_}, \texttt{val\_} and \texttt{val\_cal\_} prefixes are
  omitted. Values marked $\star$ are derived once per epoch from the logged
  pair rather than averaged over batches.}
  \label{tab:metrics-supervised}
  \begin{tabular}{@{}>{\raggedright\arraybackslash}p{0.34\textwidth}>{\raggedright\arraybackslash}p{0.60\textwidth}@{}}
    \toprule
    \textbf{Metric} & \textbf{What it measures} \\
    \midrule
    \multicolumn{2}{@{}l}{\emph{Headline}} \\
    \met{loss} & cap-scaled MSE, Eq.~\eqref{eq:capscaled}; \met{val_loss} ranks checkpoints \\
    \met{lr}, \met{epoch_time_s} & learning rate in effect; wall-clock cost of the epoch \\
    \addlinespace
    \multicolumn{2}{@{}l}{\emph{Is it learning anything?}} \\
    \met{baseline_zero_loss} & the loss of always predicting zero \\
    \met{mse_dim_i}, \met{baseline_dim_i} & per-component error and zero-predictor energy \\
    \met{r2_dim_i}\,$\star$ & $1 - \text{mse}/\text{baseline}$: 0 is no better than zero, 1 is perfect \\
    \met{persist_loss}, \met{persist_dim_i} & error of the velocity-echo predictor \\
    \met{skill_dim_i}\,$\star$ & $1 - \text{mse}/\text{persist}$: positive only when the model beats the echo \\
    \addlinespace
    \multicolumn{2}{@{}l}{\emph{What behaviour is it reproducing?}} \\
    \met{amplitude_dim_i}\,$\star$ & ratio of predicted to demonstrated spread; $<1$ is regression to the mean \\
    \met{restart_frames}, \met{restart_ratio_dim_i}\,$\star$ & how large the commanded departure is, on frames leaving standstill \\
    \met{mse_still}, \met{mse_moving} & error on hold-still targets against error on manoeuvres \\
    \met{r2_m<id>_dim_i}\,$\star$ & the per-component $R^2$ for each robot model in the fleet \\
    \addlinespace
    \multicolumn{2}{@{}l}{\emph{Is the machinery healthy?}} \\
    \met{encoder/std}, \met{encoder/mean}, \met{encoder/norm} & across-sample spread of the point-cloud descriptor; near zero means a dead encoder \\
    \met{state/cell_<k>/mean,std,max} & hidden-state norms per recurrent cell \\
    \met{grad/total}, \met{grad/<module>}, \met{grad/geometry_head} & gradient norm: globally, per sub-module, and for the auxiliary head \\
    \met{geometry_loss} & error of the auxiliary head on the geometry the frame's own cloud shows \\
    \met{geometry_explained_dim_i}\,$\star$ & the share of that geometry's spread the head accounts for; 0 is a head predicting the average \\
    \bottomrule
  \end{tabular}
\end{table}

\subsection{Reinforcement Learning Metrics}
\label{sec:metrics-rl}

The reinforcement trainer takes the network the imitation stage produced and
refines it by flying: same architecture, same weights, now shaped by the
consequences of its own actions rather than by a pilot's example. The
algorithm is a recurrent variant of Proximal Policy Optimization
\cite{schulman2017ppo} with Generalised Advantage Estimation
\cite{schulman2016gae}, described in Section~\ref{sec:rppo-gae}. Its metrics
are logged once per iteration and fall into six groups: what the policy
achieved on the task; how the policy update behaved; how well the critic is
doing; what the policy's action distribution looks like; what the reward was
made of and what the vehicle flew through; and what the run cost. The first
group is the one a reader cares about; the others exist because, at one point
or another during this work, a run failed in a way the first group could not
explain.

\subsubsection{The Six Groups}
\label{sec:metrics-rl-groups}

The columns are logged once per iteration and fall into six groups, summarised
here and defined one by one in Appendix~\ref{app:metrics}.

\paragraph{Task performance.}
What the policy achieved: the episodic return and length, the share of episodes
that ended for each reason the environment declares, the metres actually closed
between an episode's first observation and its last, and the state of the
episodes still open at the end of the rollout. This is the group a reader cares
about, and within it the end reasons and the progress closed --- not the return
--- are the primary reading, for the structural reason given below.

\paragraph{Policy optimisation diagnostics.}
How the update behaved: the two loss terms, the approximate divergence that
governs the trust region of Section~\ref{sec:rl-schedules} together with the
number of epochs actually run, the fraction of samples reaching the clipping
boundary, and the norms of the gradients before clipping. They describe the
optimiser rather than the task, and they are where a run that stops improving is
diagnosed.

\paragraph{Value function.}
Everything in the advantage rests on the critic, so its regression error, the
fraction of the return's variance it explains, and the statistics of the values
and of the advantages are logged separately. A critic predicting a constant
produces advantages that are the raw returns, and nothing else in the run says
so.

\paragraph{Exploration and action statistics.}
The width of the Gaussian, the entropy, the magnitude of the commanded actions
and the share of the distribution's mean lying outside the action envelope. They
say whether the policy is still exploring, and whether it is trying to command
what the platform cannot execute.

\paragraph{Reward decomposition and flight signals.}
Each term of the reward of Section~\ref{sec:rl-reward} contributes its own
column, and each signal the environment reports --- clearance, tilt, height
above ground, speed --- is logged as a worst and a mean over the rollout. The
decomposition is what makes a claim about the reward a number rather than an
inference: it is possible to read which term the policy is actually responding
to, and how close the fleet flew to the threshold behind each barrier.

\paragraph{Efficiency and thermal guard.}
Throughput, where the wall-clock time of an iteration went, and the time lost to
the thermal guard of Section~\ref{sec:sl-optimisation}. They do not describe the
policy, and they are what makes two runs comparable in cost.

\medskip

The first group is the one that answers the question the work asks; the other
five exist because, at one point or another during this work, a run failed in a
way the first could not explain.

\subsubsection{Summary}

Table~\ref{tab:metrics-rl} lists the reinforcement learning metrics by group.

\begin{longtable}{@{}>{\raggedright\arraybackslash}p{0.40\textwidth}>{\raggedright\arraybackslash}p{0.54\textwidth}@{}}
  \caption{Reinforcement learning metrics, one row per iteration, grouped by
  the question they answer.}
  \label{tab:metrics-rl} \\
  \toprule
  \textbf{Metric} & \textbf{What it measures} \\
  \midrule
  \endfirsthead
  \multicolumn{2}{@{}l}{\emph{Table~\thetable\ (continued)}} \\
  \toprule
  \textbf{Metric} & \textbf{What it measures} \\
  \midrule
  \endhead
  \bottomrule
  \endlastfoot
    \multicolumn{2}{@{}l}{\emph{Task performance}} \\
    \met{ep_return_{mean,std,min,max}} & episodic return of the episodes closed this iteration \\
    \met{ep_length_{mean,std}} & their length in steps \\
    \met{n_episodes}, \met{episodes_total} & episodes closed now, and since the start \\
    \met{success_rate}, \met{term/<reason>} & share of episodes that arrived, and share per end reason \\
    \met{ep_progress_{mean,std}} & metres closed per episode; the return's honest counterpart \\
    \met{inflight_{steps,budget,d_goal}} & state of the episodes still open \\
    \addlinespace
    \multicolumn{2}{@{}l}{\emph{Policy update}} \\
    \met{pg_loss}, \met{entropy}, \met{actor_loss} & clipped surrogate, policy entropy, total actor objective \\
    \met{approx_kl}, \met{approx_kl_max} & policy movement: epoch mean, and the largest mini-batch \\
    \met{target_kl}, \met{kl_early_stop}, \met{epochs_run}, \met{n_epochs} & the trust region and whether it cut the update short \\
    \met{clip_frac}, \met{clip_eps} & share of samples on the hinge; the $\epsilon$ in effect \\
    \met{replay_logratio_max} & recurrent-state correctness check; must be $\approx 0$ \\
    \met{lr_actor}, \met{lr_critic}, \met{actor_grad_norm} & optimiser state \\
    \met{critic_warmup} & 1 while only the critic is learning; the actor's columns are then flat by construction \\
    \addlinespace
    \multicolumn{2}{@{}l}{\emph{Value function}} \\
    \met{vf_loss}, \met{critic_loss}, \met{critic_grad_norm} & critic regression loss and gradient \\
    \met{explained_var}, \met{explained_var_batch} & variance of the returns the critic accounts for; rollout-level and per-batch \\
    \met{value_{mean,std,min,max}}, \met{value_mse} & distribution of $V(s)$ and its error; \met{value_std} $= 0$ is a constant critic \\
    \met{return_{mean,std}}, \met{adv_{mean,std}} & GAE returns and raw advantages \\
    \addlinespace
    \multicolumn{2}{@{}l}{\emph{Action distribution}} \\
    \met{action_std_mean} & width of the Gaussian: exploration \\
    \met{action_abs_mean}, \met{action_sat_frac} & size of the samples; how often the clamp fired \\
    \met{action_mu_abs_mean}, \met{action_mu_outside_frac} & size of the mean; how often it lies outside the envelope \\
    \addlinespace
    \multicolumn{2}{@{}l}{\emph{Reward and environment}} \\
    \met{reward_step_{mean,std}} & the reward per step \\
    \met{reward/<term>} & mean per-step contribution of each term; sums to the above \\
    \met{signal/<name>_{worst,mean,coverage}} & what the fleet flew through \\
    \met{signal/<name>_{total,worst_slice}} & growth of the slices' own counters \\
    \addlinespace
    \multicolumn{2}{@{}l}{\emph{Health and cost}} \\
    \met{obs_rms_var_mean}, \met{ret_rms_var} & scales of the observation and return normalisers \\
    \met{encoder/*}, \met{state/cell_*} & encoder spread and hidden-state norms, over the rollout \\
    \met{aux/loss}, \met{aux/explained_dim_i} & the auxiliary geometry task, when it is switched on \\
    \met{sps}, \met{*_time_s} & throughput and where an iteration's time goes \\
    \met{gpu_temp_c}, \met{cpu_temp_c}, \met{thermal_pause_s} & temperatures and time lost to the thermal guard \\
\end{longtable}
